\section{Experimental Results: Imputation Evaluation}
\label{sec:imputation_results}

We first evaluate the imputation efficacy of the diverse architectures in isolated phases. All results reflect testing on the MIMIC-IV time-series subset at a random missingness mask probability of $p=0.3$.

\subsection{Phase 1: SAITS vs. SSSD on Core Features (9-Feature Subset)}
Table \ref{tab:results_9feat} outlines performance on a restricted 9-feature subset, allowing direct contrast against SSSD's diffusion overheads.

\begin{table}[h]
\centering
\caption{Comparison on MIMIC-IV Subset (9 Features, 1,000 samples, $p=0.3$)}
\label{tab:results_9feat}
\begin{tabular}{@{}lccc@{}}
\toprule
\textbf{Model} & \textbf{MAE} & \textbf{RMSE} & \textbf{Latency (s/sample)} \\ \midrule
Linear Interp.  & 0.0337       & 0.1460        & $<0.001$                   \\ \midrule
\textbf{SAITS}  & \textbf{0.0295} & \textbf{0.3005} & \textbf{0.024}            \\
SSSD            & 0.0873       & 0.5609        & 3.690                      \\ \bottomrule
\end{tabular}
\end{table}
SAITS significantly outperforms SSSD in accuracy while generating a two-orders-of-magnitude reduction in latency. In larger features (17 variables), SSSD training crashed altogether, failing to recover standard interpolation baselines, whereas SAITS remained highly performant.

\subsection{Phase 2: Consolidation and KGI Exploration (17-Feature Paradigm)}
Expanding to 17 variables allowed robust study of graph-driven priors. Initially, static UMLS configurations caused slight performance degradation. The deployment of the SapBERT embeddings coupled to graph learning and the Parallel Gated iteration offered pronounced reductions in error compared to Random-prior networks. 

\begin{table}[h]
\centering
\caption{Final Consolidated Benchmark (17 Features, Parallel Architecture)}
\label{tab:final_benchmark}
\resizebox{\textwidth}{!}{
\begin{tabular}{@{}lccccc@{}}
\toprule
\textbf{Model} & \textbf{MAE} $\downarrow$ & \textbf{RMSE} $\downarrow$ & \textbf{MRE} $\downarrow$ & \textbf{R$^2$ Score} $\uparrow$ & \textbf{Corr. Err} $\downarrow$ \\ \midrule
Linear Interp.       & 0.4038 & 0.6145 & 1.6754 & 0.6195 & 0.0269 \\ \midrule
\textit{Deep Models (True Baselines)} \\
True Vanilla SAITS (No Graph) & 0.4566 & 0.6526 & 1.7209 & 0.5742 & 0.0386 \\
\textbf{True Prior Nullo (Random Graph)} & \textbf{0.4025} & \textbf{0.5998} & \textbf{1.5713} & \textbf{0.6402} & \textbf{0.0483} \\ \midrule
\textit{Deep Models (Parallel Gated)} \\
SapBERT+CI-GNN (Par)          & 0.4028 & 0.5984 & 1.5686 & 0.6420 & 0.0520 \\ \bottomrule
\end{tabular}
}
\end{table}

Interestingly, an ablation architecture---the \textbf{True Prior Nullo} (randomly initialized, no formal UMLS prior penalty graph loss)---operated exceptionally well (MAE 0.4025). Analysis of the $A_{learn}$ graph converged upon by True Prior Nullo highlighted that 100\% of the strongest data-driven attention edges corresponded to authentic UMLS priors (Precision=1.00), implicitly discovering links between Systolic/Diastolic BP ($A=0.74$) without explicit graph supervision. This hints that strict static UMLS priors are redundant when network capacity and semantics are robustly initialized.

\subsection{Phase 3: Deep Baselines and Foundational Models on SOTA Data Framework}
Integrating additional rigorous standards, we tested deeper baselines alongside TimesFM, achieving a holistic SOTA assessment snapshot in our heavily correlated configuration.

\begin{table}[h]
\centering
\begin{tabular}{lcccccc}
\toprule
\textbf{Model} & \textbf{Epoch} & \textbf{MAE} $\downarrow$ & \textbf{RMSE} $\downarrow$ & \textbf{MRE} $\downarrow$ & \textbf{R2} $\uparrow$ & \textbf{Corr. Err} $\downarrow$ \\
\midrule
\textbf{SAITS (Vanilla)} & 48 & \textbf{0.3245} & \textbf{0.5092} & \textbf{1.6735} & \textbf{0.7392} & 0.1109 \\
SAITS (SapBERT Embeds) & 49 & 0.4462 & 0.6300 & 2.2403 & 0.6008 & 0.0922 \\
BRITS & 49 & 0.5277 & 0.7235 & 1.8381 & 0.4739 & 0.1585 \\
TimesFM 2.5 (Zero-shot) & - & 0.4524 & 0.6960 & 1.9955 & 0.4899 & \textbf{0.0405} \\
\midrule
\textbf{KGI-SAITS (Generic, No-GNN)} & \textbf{70} & \textbf{0.3381} & \textbf{0.5230} & \textbf{1.7059} & \textbf{0.7248} & \textbf{-} \\
\bottomrule
\end{tabular}
\caption{Comprehensive Evaluation on SOTA Dataset. SAITS achieves high point-wise accuracy, while TimesFM Foundation Models capture robust correlation structure.}
\label{tab:results_sota}
\end{table}

While SAITS Vanilla reached prime reconstruction accuracies ($\sim$0.3245 MAE), the Foundation Model \textbf{TimesFM 2.5} drastically minimized correlation error (0.0405 vs. 0.1109) strictly through zero-shot global temporal modeling semantics, showcasing an alternative perspective for preserving clinical covariance structures.

\subsection{Phase 4: Impact of High-Fidelity Relational Verbalization}
In an attempt to enrich the knowledge graph injection, we engineered a ``High-Fidelity'' variant of KGI-SAITS. Instead of generic temporal alignment prompts (e.g., ``is associated with''), this variant utilized explicit hierarchical clinical relationship templates derived from UMLS (e.g., ``Heart Rate is a type of Vital sign'').

Counterintuitively, this High-Fidelity approach \textbf{deteriorated} the imputation accuracy. Evaluated on the SOTA dataset under identical conditions, the High-Fidelity KGI model yielded an MAE of 0.4935 and $R^2$ of 0.5431, trailing behind both the Generic KGI (MAE 0.4575, $R^2$ 0.5886) and the SAITS Vanilla baseline (MAE 0.4596, $R^2$ 0.5888). 

We hypothesize that integrating complex, precise multi-hop clinical syntax produces MedBERT embeddings that are excessively intricate for the continuous continuous-time forecasting domain. The model's cross-attention mechanisms likely become saturated by parsing lengthy linguistic structures rather than focusing on the fundamental topographical proximity of the features, indicating that generic, standardized templating serves as a more effective medium for graph-topology translation in time-series modules.
